\section{Does the Story Continue? A Covariant Unification in 12d}

In the previous section we saw that the D(-1) and D7 solutions admit an uplift to 12d, and that their Killing spinor equations are reproduced by reducing a 12d covariantly-constant spinor condition on a torus, were such a 12d theory to exist.
It is natural to ask whether other brane-like solutions in the type IIB theory have an analogous 12d uplift.


\subsection{Brane-like solutions}
\label{sec:general brane solns}
It has been shown \cite{Stelle:1998xg} that, working in the Einstein frame, with $F_n$ an n-form and the action
\begin{equation}
S = \int d^Dx\sqrt{-g}\left[
    R - \frac{1}{2}(\partial\phi)^2 - \frac{1}{2}e^{\alpha\phi}|F_n|^2
\right],
\end{equation}
the equations of motion, together with Bianchi identity (assuming $F = dA$) are given by
\begin{equation}
\begin{aligned}
R_{MN}
&= \frac{1}{2}\left(
    \partial_M\phi\partial_N\phi + e^{\alpha\phi}S_{MN}\big|_F
\right)\\
dF&=0\\
d[e^{\alpha\phi}*_D F] &=0\\
\Box\phi&=\frac{\alpha}{2}e^{\alpha\phi}|F|^2
\end{aligned}
\end{equation}
where
\begin{equation}
S_{MN}\big|_{F}\equiv
\frac{1}{2}\left[
    \frac{1}{(n-1)!}F_{M...}F_N^{...}
    -\frac{n-1}{D-2}\left(\frac{1}{n!}F^2\right)g_{MN}
\right],
\end{equation}
and $\Delta\equiv \alpha^2 +\frac{2 d \tilde{d}}{D-2}$.
There exists the electric solution, defined as the case where the form-field potential has legs along the worldvolume directions,
\begin{equation}
\begin{aligned}
ds_D^2
&=H^{-\frac{4\tilde d}{\Delta(D-2)}}ds_{d}^2+
H^{\frac{4d}{\Delta(D-2)}}ds_{D-d}^2,\quad
d = n-1\\
F_{\mu_0\mu_1...\mu_{n-2},m}
&=\frac{2}{\sqrt\Delta}\varepsilon_{\mu_0\mu_1...\mu_{n-2}}\partial_m (H^{-1})\\
e^\phi
&=H^{2\alpha/\Delta}
\end{aligned}
\end{equation}
as well as the magnetic solution, defined as the case where the field strength has its legs along the transverse directions,
\begin{equation}
\begin{aligned}
ds_D^2
&=H^{-\frac{4\tilde d}{\Delta(D-2)}}ds_{d}^2+
H^{\frac{4d}{\Delta(D-2)}}ds_{D-d}^2,\quad
d = D-n-1\\
F_{m_1...m_n}
&= \frac{2}{\sqrt{\Delta}}\varepsilon_{m_1...m_n m}\partial_m H\\
e^\phi &= H^{-2\alpha/\Delta}
\end{aligned}
\end{equation}
In 11d supergravity, type IIA, and type IIB, one finds $\Delta = 4$ for all branes.
We provide a catalog of the various brane solutions in Appendix \ref{sec:catalog-of-branes}.

Disregarding the fact that no 12d supergravity theory is known to exist, one can immediately identify the difficulties in finding 12d uplifts of type IIB branes.

One can readily see that there does not exist nontrivial ways to embed type IIB metric, scalars, and form fields into 12d metric and form fields.
At the level of supergravity, 12d fields transform under the little group SO(10) while the 10d fields transform under SO(8).
In Table~\ref{tab:10d-12d-dof} we summarize the 10d massless on-shell degrees of freedom in type IIB, alongside the on-shell degrees of freedom of the corresponding 12d field of the same form rank.

\begin{table}[h]
\centering
\begin{tabular}{l|ccccccc}
\hline
 & metric & scalar & 1-form & 2-form & 3-form & 4-form & 5-form \\
\hline
Type IIB $SO(8)$ & $35$ & $2\times 1$ & $0\times 8$ & $2\times 28$ & $0\times 56$ & $1\times 35$ & $-$ \\
12d $SO(10)$ & $54$ & $-$ & $10$ & $45$ & $120$ & $210$ & $126$ \\
\hline
\end{tabular}
\caption{
On-shell bosonic degrees of freedom of the type IIB fields, given as (number of fields) $\times$ (dimension of the $SO(8)$ little-group representation), compared with the dimension of the $SO(10)$ representation of a 12d field of the same form rank.
}
\label{tab:10d-12d-dof}
\end{table}

At the level of fields, then, a fully 12d uplift is not available: no configuration of a 12d graviton and 12d form fields reduces to the type IIB field content, and we should not expect type IIB to arise as the dimensional reduction of a genuine 12d (super)gravity realizing the full 12d spacetime symmetries.

This does not, however, exhaust the 12d interpretation.
Already in the axio-dilaton sector the embedding into the 12d metric was only partial---the two off-diagonal Kaluza-Klein 1-forms were set to zero, so the full 12d metric was never reconstructed---yet it still furnished a genuine 12d perspective on the D(-1) and D7 branes, as discussed in the previous section.
In the same spirit, in the following subsection we organize the type IIB fields into 12d objects \emph{without} realizing the full 12d spacetime symmetries: not all degrees of freedom of the 12d gravity and form fields are excited, and the construction is best understood as a covariant repackaging of the 10d theory rather than the reduction of a true 12d one.

\subsection{A Covariant Unification in 12d}

The various brane-related evidence for an effective 12d interpretation of the type IIB theory suggests a very specific 12d interpretation of the axio-dilaton sector.
In addition, one seeks a 12d interpretation of the type IIB RR and NSNS form fields.
In this section, we gather the various 12d insights obtained in the previous section from the type IIB branes, and present a 12d covariant unification of $SL(2,R)\times SO(9,1)$ form fields.
We will use hats $\hat{\;}$ to denote fields in the higher dimension.
The 10d metric $g_{mn}$ with $m,n=0,1,...,9$ will be interpreted as embedded inside a 12d metric $\hat{g}_{MN}$, with $M,N=0,1,...,11$. We will use the calligraphic $\hat R$ to denote the Ricci scalar computed with $\hat{g}_{MN}$, and use $\hat{F}_{n+1} = d\hat{C}_{n}$ to denote form fields in 12d.

Let the 12d coordinates be parameterized by $(x^m,z_1,z_2)$, and let $\hat{g}_{ab}$ be the $2\times 2$ metric on the torus.
The key insight from the previous section is that the axio-dilaton sector of type IIB may arise from the 12d metric embedding given by
\begin{equation}
\hat{g}_{MN}=\begin{pmatrix}
g_{mn} & 0 \\
0 &\hat{g}_{ab}
\end{pmatrix},\quad
\hat{g}_{ab}
=\frac{1}{\tau_2}\begin{pmatrix}
1& \tau_1\\
\tau_1 & \tau_1^2+\tau_2^2
\end{pmatrix}.
\label{12d metric ansatz}
\end{equation}
In particular,
\begin{equation}
\begin{aligned}
\frac{1}{2\kappa^2}\int d^{10}x\sqrt{-g}\left(
R - \frac{|\nabla \tau|^2}{2\tau_2^2}
\right)
&=
\frac{Vol(T_2)}{2\kappa_{12}^2}\int d^{10}x \sqrt{-\hat{g}} \hat{R}\\
&=\frac{1}{2\kappa_{12}^2}\int_{T_2}dz_1 dz_2\int d^{10}x\sqrt{-\hat{g}}\hat{R},
\end{aligned}
\label{axio-dilaton 12d uplift}
\end{equation}
where we have schematically defined $\kappa_{12}$ by 
\begin{equation}
\frac{1}{\kappa^2} = \frac{Vol(T_2)}{\kappa_{12}^2}
=\frac{1}{\kappa_{12}^2}\int_{T_2}*_2 1,\quad
[\kappa_{12}^2] = L^{10}.
\end{equation}
The 3-form field strengths form an SL(2,R) doublet. 
To unify them in 12d we define a 12d 4-form field strength with exactly one leg on the torus:
\begin{equation}
\hat{F}_4=d\hat C_3,\quad
\hat C_3 \equiv B_2\wedge dz_1 + C_2\wedge dz_2.
\end{equation}
Using the 12d metric ansatz \eqref{12d metric ansatz}, the 10d 3-form field strengths and their axio-dilaton couplings follow from contracting $\hat F_4$:
\begin{equation}
\begin{aligned}
|\hat{F}_4|\bigg|_{\hat{g}_{MN}}
&=\hat{g}^{z_1 z_1}|H_3|^2+2\hat{g}^{z_1 z_2}|F_3\cdot H_3|+\hat{g}^{z_2 z_2}|F_3|^2\\
&=\left(
e^{-\Phi}|H_3|^2 + e^\Phi|F_3-CH_3|^2
\right)\bigg|_{g_{mn}}.
\end{aligned}
\end{equation}
Finding a 12d interpretation for the 5-form sector is trickier. 
One observes that in 10d, the composite SL(2, R) singlet 5-form 
\begin{equation}
\tilde F_5 = F_5 - \frac12 C_2 \wedge H_3 + \frac{1}{2}B_2\wedge F_3    
\end{equation}
can not be sensibly constructed in 12d, due to the lack of a pair of form field potential and strength with ranks that sum to 5 in 12d.
However, the 10d self-dual 5-form field strength admits two possible uplifts to 12d:
\begin{equation}
\hat{F}_5 = F_5,\quad \hat{F}_7=F_5\wedge dz_1 \wedge dz_2.
\end{equation}
Then note that
\begin{equation}
\begin{aligned}
|\hat{F}_5|^2\bigg|_{\hat{g}_{MN}} 
&= |F_5|^2\bigg|_{g_{mn}},\\
|\hat{F}_7|^2\bigg|_{\hat{g}_{MN}}
&=\det(\hat{g}_{ab})|F_5|^2\bigg|_{g_{mn}} = |F_5|^2\bigg|_{g_{mn}}.
\end{aligned}
\end{equation}
Then we may define a 12d 7-form
\begin{equation}
\tilde{\hat{F}}_7 \equiv \hat{F}_7 + \frac{1}{2}\hat{C}_3 \wedge \hat{F}_4
\end{equation}
that exactly supplies the type IIB composite 5-form contribution upon contraction in 12d:
\begin{equation}
|\tilde{\hat{F}}_7|^2\bigg|_{\hat{g}_{MN}} = 
\det(\hat{g}_{ab})|\tilde F_5|^2\bigg|_{g_{mn}}.
\end{equation}
The additional utility of $\hat{F}_7$ is that it allows us to write the 10d self-duality condition on $F_5$ as a 12d Hodge duality condition
\begin{equation}
\hat{F}_7 = *_{12}\hat{F}_5\quad
\Leftrightarrow\quad
F_5 = *_{10}F_5.
\label{12d interpretation of 10d self-duality}
\end{equation}
The 10d Chern-Simons term can also be obtained using the 12d potentials and their corresponding field strengths:
\begin{equation}
\frac{1}{2}\int_{T_2\times \mathbb{R}^{1,9}} \hat{C}_4\wedge \hat{F}_4\wedge\hat{F}_4=
Vol(T_2)\int_{\mathbb{R}^{1,9}}C_4\wedge H_3\wedge F_3.
\end{equation}
We note that reproducing the type IIB Chern-Simons term in 12d necessitates the inclusion of both 4- and 5-form field strengths.
One can write down a ``12d"\footnote{
Not dynamical 12d, but dynamical 10d times a 2-torus.
}
action
\begin{equation}
\begin{aligned}
S_{IIB}=
S_{\text{``12"}}
= \frac{1}{2\kappa_{12}^2}\int_{T_2}dz_1 dz_2\int d^{10}x
\sqrt{-\hat{g}}
\left(
\hat{R} - \frac{1}{2}|\hat F_4|^2-\frac{1}{4}|\tilde{\hat{F}}_7|^2
\right)-\frac{1}{4\kappa_{12}^2}\int_{T_2\times \mathbb{R}^{1,9}}\hat{C}_4\wedge \hat{F}_4\wedge \hat{F}_4.
\end{aligned}
\label{12d action}
\end{equation}
The $\hat F_7$ does not arise from an independent degree of freedom, it is related to $\hat C_4$ by $*_{12}d\hat C_4 = \hat F_7$.
The action \eqref{12d action} is exactly the Einstein frame type IIB action \eqref{Einstein frame IIB action} with $g_s=1$.
The 2d integrand is just a repackaging of $1/\kappa^2$, and it is likely that $\hat C_3, \hat C_4$ do not furnish independent degrees of freedom, as has been discussed in \cite{Tseytlin:1996ne}.

\paragraph{The Volume of the Torus}

\subsection{Uplift of type IIB branes}
We now present the 12d "uplifts" of the various type IIB brane solutions, using the effective 12d "action".


We would like to ask the same question about type IIB.
Despite one often think of F1 and NS5, D1 and D5 as EM dual pairs, their form fields $H_3$ and $F_3$ configuration are not related by EM duality, in 10d, and of course they have opposite dilaton profiles as well\footnote{The brane solutions are listed in \ref{sec:10d IIB branes}}.

To obtain the NS5 from F1 by EM duality, one actually has to perform an accompanying SL(2, Z) transformation to invert the dilaton.
Similar situation is encountered on the D3 brane worldvolume \cite{Tseytlin:1996it}.
Recently a 12d uplift of the various form fields have been proposed \cite{Cheng:2025efp}. Let us uplift the F1, D1, NS5, D5 branes to 12d.

\paragraph{F1}
\begin{equation}
\begin{aligned}
ds^2
&= H^{-3/4}ds_{1,1}^2
+ H^{1/4}ds_8^2
+H^{-1/2}dz_1^2
+H^{1/2}dz_2^2\\
(\mathcal{F}_4)_{\mu_0 \mu_1 m z_1}
&= \varepsilon_{\mu_0 \mu_1}\partial_m H^{-1}
\end{aligned}
\end{equation}

\paragraph{D1}
\begin{equation}
\begin{aligned}
ds^2
&= H^{-3/4}ds_{1,1}^2
+ H^{1/4}ds_8^2
+H^{1/2}dz_1^2
+H^{-1/2}dz_2^2\\
(\mathcal{F}_4)_{\mu_0 \mu_1 m z_2}
&= \varepsilon_{\mu_0 \mu_1}\partial_m H^{-1}
\end{aligned}
\end{equation}


\paragraph{NS5}
\begin{equation}
\begin{aligned}
ds^2
&= H^{-1/4}ds_{1,5}^2
+ H^{3/4}ds_4^2
+H^{1/2}dz_1^2
+H^{-1/2}dz_2^2\\
(\mathcal{F}_4)_{m_1 m_2 m_3 z_1}
&= \varepsilon_{m_1 m_2 m_3 m}\partial_m H
\end{aligned}
\end{equation}

\paragraph{D5}
\begin{equation}
\begin{aligned}
ds^2
&= H^{-1/4}ds_{1,5}^2
+ H^{3/4}ds_4^2
+H^{-1/2}dz_1^2
+H^{1/2}dz_2^2\\
(\mathcal{F}_4)_{m_1 m_2 m_3 z_2}
&= \varepsilon_{m_1 m_2 m_3 m}\partial_m H
\end{aligned}
\end{equation}

\paragraph{Caveat: How ``12d" is this uplift?}
It is important to dicuss the nature of these uplifts.
In any dimension, with a Einstein-Hilbert action couled to $n$-form field strength, the Einstein equations is $R_{MN} = S_{MN}$ where $S_{MN}$ is the trace-reversed stress energy tensor of the form field,
\begin{equation}
S_{MN} = \frac{1}{2}\left(\frac{1}{(n-1)!}F_{M...}F_{N}{}^{...} - \frac{n-1}{D-2}\frac{1}{n!}F^2g_{MN}\right).
\label{trace-reversed stress-energy tensor}
\end{equation}
An important distinction between the form-field uplift ansatz and the usual Kaluza-Klein uplift (e.g. in IIA to 11d) is that not all DoFs of the higher-dimensional form field is turned on.
This has important consequences on whether one may ``interpret" the uplifted solution as being in the higher dimensional theory.
We note that, for the D1, F1, NS5, D5 uplifts above, they do NOT satisfy the 12d Einstein equation $R_{MN} = S_{MN}$.
For example, for the F1 uplift given above, one can compute the $R_{\mu\nu}, R_{zz}, R_{mn}$, with the uplifted 12d metrics, but in order for say $R_{\mu\nu}$ to equare $S_{\mu\nu}$, one needs to compute $S_{\mu\nu}$ with $\mathcal{F}_4$ with $n=3$ instead of $n=4$, as well as $D=10$ instead of $D=12$, in \eqref{trace-reversed stress-energy tensor}.
This is because the factors related to $n!, (n-1)!, D$ in the formula for the trace-reversed stress-energy tensor \eqref{trace-reversed stress-energy tensor} are present to discount the overcounting in trace-reversal as well as contracting the form field. But with the ansatze $\mathcal{F}_4 = F_3 \wedge dz_1$, $\mathcal{F}_4$ fundamentally has the DoFs of a 10d 3-form.
So the notion of 12d form-field is much weaker than the notion of 12d gravity.
Rather, the ``12d" uplift is more appropriately described as a repackaging of the 10d equations of motion.
Again, this all have to do with the fact that not all DoFs of 12d form fields have a place in the 10d theory.

\paragraph{Interpreting the EM duality of Uplift}

In 12d the SL(2, Z) transformation $e^\Phi \to e^{-\Phi}$ becomes $z_1 \to z_2, z_2 \to -z_1$, which indeed maps between the F1 and D1 and between the NS5 and D5 solutions.
There is a pattern here: the F1 uplift's form field configurations and that of the D5 uplift are very close to being related by EM duality in 12d, and so are the D1 and NS5 uplifts.

Let us perform EM duality on the $\mathcal{F}_4$ configuration of F1.
\begin{equation}
(\mathcal{F}_8)_{m_1 m_2 ...m_7 z_2} = \varepsilon_{m_1 m_2 m_3 ... m_7 n}\partial_n H
\end{equation}
if we decompose $\mathcal{F}_8$ back to $\mathcal{F}_4$ in 12d, by $\mathcal{F}_8 = \varepsilon_4 \wedge \mathcal{F}_4$, this is indeed the $\mathcal{F}_4$ configuration of the D5 uplift. From there performing SL(2, Z) duality $z_1 \leftrightarrow z_2$ we reach the NS5 uplift.

\paragraph{Comments on Role of dilaton and F-theory}
Usually, in supergravity theories, the dilaton is the volume of the compact manifold.
In type IIB, however, the compact torus has fixed volume: as the radius of one circle grows, the other shrinks.
The natural place for the non-compact scalar dilaton to live is the radius of one of the circles $R_1 \sim e^{\Phi}$, while the radius of the other circle shall scale inversely $R_2 \sim e^{-\Phi}$.
In this way, type IIB backgrounds can not have a 12d decompactifying limit, because taking the decompactifying limit of one circle immediately compactifies the other circle. 
The exception being the D(-1)-instanton uplift, where the pp-wave propages inside a non-comapct torus.
The KK-Monopole uplift of D7-brane is not an exception because a 12d KK-monopole is localised on $S_1 \times R^3$: there is indeed a compact circle in 12d.

While the two circle radius shall scale inversely to each other, one may consider the case where they are equal: $e^\Phi = e^{-\Phi} = 1$.
This implies that the radius of both circles are finite hence must be a compactification limit of the full torus.
